\chapter{Tests et Sécurité} 
\minitoc
\newpage

\section{Introduction}
Tout produit logiciel professionnel, avant d'être mis entre les mains de ses utilisateurs finaux, doit passer par une phase rigoureuse de qualification et d'Assurance Qualité (QA). Ce chapitre détaille la stratégie de test mise en œuvre pour garantir l'absence de régressions, ainsi que les mesures de sécurité préventives appliquées sur l'architecture backend.

\section{Stratégie de Test et Assurance Qualité}
La complexité inhérente à une application gérant des programmes d'entraînement asynchrones et des relations inter-utilisateurs multiples rend les tests manuels insuffisants.

\subsection{Tests d'intégration avec Postman}
Pour valider le fonctionnement de notre API REST (le cœur du système), nous avons utilisé Postman de manière systématique.
Plutôt que de tester de manière ad hoc, nous avons créé une \textbf{Postman Collection} automatisée. Chaque scénario critique a été associé à des scripts de test (écrits en JavaScript au sein de Postman) permettant de valider :
\begin{itemize}
    \item \textbf{Le code HTTP de retour :} S'assurer qu'une requête valide retourne `201 Created` ou `200 OK`, et qu'une tentative d'accès sans token JWT retourne bien `401 Unauthorized`.
    \item \textbf{Le schéma JSON :} Vérifier que l'API renvoie bien les champs attendus (ex: `program.id`, `program.days`).
    \item \textbf{Les contraintes de temps de réponse :} S'assurer que le calcul du Dashboard nutritionnel s'effectue en moins de 300 ms.
\end{itemize}

\subsection{Tests fonctionnels sur Flutter}
Côté client, le développement Flutter s'est accompagné d'une phase de tests intensifs sur divers émulateurs (Pixel 6, iPhone 14 Pro) et sur un périphérique physique.
L'objectif était de tester la résilience de l'application face à des situations dégradées :
\begin{itemize}
    \item \textbf{Perte de connexion réseau :} Vérification de l'affichage des messages d'erreur "Offline" sans crash de l'application.
    \item \textbf{Gestion des états asynchrones :} S'assurer que les indicateurs de chargement (CircularProgressIndicator) bloquent correctement les boutons pour éviter les doubles requêtes POST de la part de l'utilisateur.
\end{itemize}

\section{Sécurité de l'Application}
La manipulation de données relatives à la santé, aux mensurations et à la messagerie privée exige un niveau de sécurité optimal. Plusieurs couches de sécurité ont été implémentées.

\subsection{Sécurité de l'authentification et Hachage}
Comme abordé lors du Sprint 1, aucun mot de passe n'est stocké en clair. Le recours à \texttt{bcrypt} (avec un coût de salage à 10) rend les attaques par dictionnaire (Rainbow Tables) caduques.
La transmission de l'identifiant se fait systématiquement encapsulée dans un \textbf{Token JWT}. Ce jeton est doté d'une date d'expiration (TTL - Time To Live) courte. 

\subsection{Protection contre les attaques courantes (OWASP)}
L'API Node.js a été renforcée contre les vulnérabilités les plus fréquentes recensées par l'OWASP (Open Web Application Security Project) :
\begin{itemize}
    \item \textbf{Injections SQL :} L'utilisation exclusive de l'ORM TypeORM prévient nativement ce type d'attaque grâce à la paramétrisation automatique des requêtes. Aucune requête brute basée sur de la concaténation de chaînes utilisateur n'est exécutée.
    \item \textbf{CORS (Cross-Origin Resource Sharing) :} Un middleware restrictif a été configuré pour empêcher des domaines non autorisés de consommer l'API depuis un navigateur web. Seul le Dashboard Angular officiel est "whitelisté".
    \item \textbf{Rate Limiting :} Afin de se prémunir contre les attaques par force brute (Brute Force) sur le flux de connexion, un middleware \texttt{express-rate-limit} bloque temporairement l'adresse IP source après plusieurs tentatives de login infructueuses.
\end{itemize}

\section{Conclusion}
Ce huitième chapitre a démontré que notre projet n'est pas qu'une simple maquette académique, mais une solution logicielle répondant aux normes de l'industrie professionnelle. La stratégie de test assure la fiabilité des flux complexes et les mesures de sécurité garantissent la confidentialité des utilisateurs. L'ensemble de ces efforts de qualification vient sceller l'achèvement technique de notre projet de fin d'études.
